\clearpage
\SourceLeaf{362}%
\pdfbookmark[0]{Notices. Mollweide: Archimedes. 1808.}{anzeigen-mollweide-1808}
\begin{center}
\rule{.85\textwidth}{.3pt}\par
Göttingische gelehrte Anzeigen.\quad 9 January 1808.\par
\rule{.85\textwidth}{.3pt}
\end{center}
\hspace*{\parindent}As is well known, in his \textit{Circuli dimensio} \textsc{Archimedes} determined bounds for the circumference of a circle by enclosing it between the perimeters of a circumscribed and an inscribed 96-gon. He calculates these numbers, or more precisely finds a number greater than the former and one smaller than the latter, by proceeding in stages from the hexagon to the 12-gon, from that to the 24-gon, and so on. For both 96-gons he therefore starts, in our terms, from an approximate value of the irrational quantity \(\BareRoot{3}\). One of these values, namely \(\frac{265}{153}\), is slightly too small, and the other, \(\frac{1351}{780}\). is slightly too large. The former is used for the circumscribed polygons, the latter for the inscribed ones. On closer inspection these approximate values are found in the sequence \(\frac21,\frac53,\frac74,\frac{19}{11},\frac{26}{15}\) and so on, whose terms are alternately greater and smaller than \(\BareRoot{3}\), and each of which differs less from it than any other fraction expressed by smaller numbers. The fraction \(\frac{265}{153}\) is the eighth term of the sequence and \(\frac{1351}{780}\) the eleventh. It therefore seems that \textsc{Archimedes} found these approximate values by a method rather than by chance. Yet he gives no account at all of how he arrived at them, and there is otherwise no indication that the ancients knew our methods of solving such problems. This leaves room for conjecture. Professor \textsc{Mollweide} of Halle has investigated the matter in a short memoir recently sent to the Royal Society, of which he is a correspondent, entitled
\begin{center}\textit{De methodo ab Archimede adhibita ad rationem, in qua inter se sunt latus trianguli aequilateri et radius circuli circumscripti, numeris veritati proxime exprimendam}\end{center}
\noindent He gives a procedure suited to the state of arithmetic among the ancients, which may therefore be the very one \textsc{Archimedes} used. Mr. M. denotes the side of the triangle by \(AC\) and the radius of the %
